Take O as origin and let $|A|=|B|=|C|=8$. The circumcentre D of OAB satisfies $D\cdot A=D\cdot B=32$. Thus D is the intersection of the tangents to the radius-four circle centred at O at $\dfrac{A}{2}$ and $\dfrac{B}{2}$. Similarly E and F are the other pairwise tangent intersections. Since ABC is acute, O lies inside it, and these three tangents enclose O. Therefore that radius-four circle is the incircle of DEF, giving $\boxed4$. The two given side lengths are unnecessary.
\begin{asy}
import graph;import geometry;usepackage("amsmath");size(7cm);
pair A=8*dir(90),B=8*dir(210),C=8*dir(330),D=8*dir(150),E=8*dir(270),F=8*dir(30);draw(A--B--C--cycle,gray);draw(D--E--F--cycle,blue);draw(circle((0,0),4));dot((0,0));label("$O$",(0,0),(1,0));label("$D$",D,(-1,0));label("$E$",E,(0,-1));label("$F$",F,(1,0));label("tangent construction (schematic)",(0,-10));
\end{asy}

